\section{Connections with non-K\"ahler geometry}
\label{sec:non-kahler-geometric-problems}

The general theory should explain geometric obstructions as well as solve equations. We describe six connections. In each case the distinction between a closed Bott--Chern class and an arbitrary smooth background is essential. The literature statements below are recorded with their hypotheses; the subsequent questions ask what the constructions of this book could add.

\subsection{Nef classes, bigness and the Fujiki class}

Let $X$ be a connected compact complex manifold of dimension $n$, fix a Hermitian form $h$ and a smooth strictly positive probability form $\mu$, and let $\alpha$ denote a closed real $(1,1)$-form or its Bott--Chern class. Integrals of closed classes are computed using smooth representatives. Use $v_+(h)$ for the bounded-potential mass supremum, $\lambda_\mu$ and $\mathcal V_\mu$ from \cref{def:gaud-variational-volume}, and $D_\mu$ from \cref{sec:gauduchon-collapse}.

A compact complex manifold belongs to the Fujiki class $\mathcal C$ if it is bimeromorphic to a compact K\"ahler manifold. The classical characterization by the existence of a closed K\"ahler current is due to Demailly--P\u aun. Their transcendental Grauert--Riemenschneider conjecture asks whether a closed nef real $(1,1)$-class $\alpha$ with
\[
 \int_X\alpha^n>0
\]
is big on an arbitrary compact complex manifold. Equivalently in this situation, the proposed conclusion is $X\in\mathcal C$; the source statements are recorded in \cref{sec:history-geometry}.

The bounded mass property is closely tied to this problem. Under exactly the preceding nefness and positive top-intersection hypotheses, Guedj--Lu prove
\[
 \alpha\text{ is big}\quad\Longleftrightarrow\quad v_+(h)<\infty
\]
for a Hermitian metric $h$. Boucksom--Guedj--Lu also characterize $\mathcal C$ by bounded mass together with positivity of their upper class volume for some closed Bott--Chern class. That upper volume is not the invariant $\mathcal V_\mu$ defined here.

The failure of bounded mass on the Hopf threefold does not refute the Demailly--P\u aun conjecture. A counterexample would also need a closed nef class with positive top self-intersection. Conversely, a proof of the conjecture cannot impose bounded mass as an unexplained general fact. The Hermitian framework must account for why that property follows in the particular geometry forced by such a class.

The elliptic threefold of \cref{thm:vanishing-lower-mass} supplies the complementary test: it has bounded mass, but no big closed class. Its entire nef cone is the nonnegative span of two smooth rank-one classes, and every top self-intersection is zero, by \cref{prop:elliptic-bott-chern-cones}. Hence this lower-mass example also leaves the geometric conjecture untouched. The exact boundary calculation in \cref{prop:elliptic-boundary-volume} describes its degeneration without identifying a nonclosed metric with a closed class.

The volume of this book separates two possible steps:
\begin{equation}\label{eq:dp-two-bridges}
 \int_X\alpha^n>0
 \quad\stackrel{?}{\Longrightarrow}\quad
 \mathcal V_\mu(\alpha)>0
 \quad\stackrel{?}{\Longrightarrow}\quad
 \alpha\text{ is big}.
\end{equation}
Neither arrow is established here in the unrestricted compact complex setting. The second is the closed special case of \cref{conj:nef-gaud-volume-big}. On a compact K\"ahler manifold both follow from the usual nef volume formula and the classical nef-and-big criterion. Without closedness, positivity of $\mathcal V_\mu$ cannot characterize $\mathcal C$: every Hermitian metric $h$ already satisfies $\mathcal V_\mu(h)>0$, on every compact complex manifold.

\begin{problem}\label{qu:closed-nef-gaud-bridge}
Let $\alpha$ be closed and nef with $\int_X\alpha^n>0$. Prove a positive lower bound for $\lambda_\mu(\alpha)$, and rule out the collapse in \cref{prop:gaud-positive-volume-collapse} for a sequence minimizing the bigness threshold of $\alpha$.
\end{problem}

This is a proposed route to the geometric conjecture, not a claimed weaker substitute already known to suffice automatically. The relevant test metrics can be normalized by $\int_Xh\wedge\gamma^{n-1}=1$. A positive lower bound for $D_\mu(\gamma)$ along one threshold-minimizing sequence would finish the second step. The Gauduchon equation can change the volume density while retaining an Aeppli class. Since $\alpha$ is closed, it also retains the numerator paired against $\alpha$. The missing estimate is control of the trace normalization and determinant denominator after that change. Merely prescribing a volume density does not give this estimate.

\subsection{Positive lower volume and concentration}

Let $X$ be a connected compact complex manifold of dimension $n$ and let $h$ be a Hermitian form. Its bounded-potential mass infimum and supremum are $v_-(h)$ and $v_+(h)$. For $n\geq2$, write $\mathscr G_h=\{T\geq0:\ T\text{ has bidegree }(n-1,n-1),\ \ddc T=0,\ \int_Xh\wedge T=1\}$. The measure $\Gamma_0(T)$ is the determinant-weighted measure of \eqref{eq:weak-gaud-functional} with logarithmic weight zero and the volume form chosen below.

Positive lower mass fails in general. The threefold in \cref{thm:vanishing-lower-mass} has $v_-(h)=0$ for every Hermitian metric, while \cref{prop:lower-mass-balanced-metric} supplies a balanced metric with finite upper mass. The Hopf threefold gives the opposite combination, $v_-(h)=6$ and $v_+(h)=\infty$ in its stated normalization. The upper and lower mass properties are independent.

For $n\geq2$, the determinant measures give an exact analytic formulation. Fix a smooth positive probability measure $\mathrm{d}V$, and write
\[
 \Gamma_0(T)=k_T\,\mathrm{d}V,\qquad T\in\mathscr G_h.
\]
By \cref{thm:gaud-determinant-endpoint},
\begin{equation}\label{eq:pvp-endpoint-geometric}
 v_-(h)>0
 \quad\Longleftrightarrow\quad
 \sup_{T\in\mathscr G_h}
       \|k_T\|_{L^{n/(n-1)}(\mathrm{d}V)}<\infty.
\end{equation}
The exact supremum is $v_-(h)^{-1/n}$, including the value $+\infty$. Uniform estimates for every smaller exponent $1<q<n/(n-1)$ hold without either volume property. The lower-mass example violates the endpoint already for smooth strictly positive tests. Thus positivity, $\ddc T=0$ and normalized trace alone do not control concentration at that exponent.

\begin{problem}\label{qu:geometric-gaud-endpoint}
Find geometric conditions that imply the endpoint estimate in \eqref{eq:pvp-endpoint-geometric}, and identify how they exclude the concentration mechanism of \cref{sec:vanishing-lower-mass}. For a fixed smooth density, determine whether the smaller family of maximizing Gauduchon currents admits stronger estimates than the full normalized cone.
\end{problem}

The balanced condition and bounded upper mass are insufficient, even together. The explicit identity \eqref{eq:lower-mass-vertical-energy} relates the failing lower bound to a vertical Dirichlet integral on an elliptic fibration. Meanwhile the horizontal curvature forms $\alpha,\beta$ of \cref{sec:vanishing-lower-mass} have a positive exact product $\alpha\wedge\beta$, which excludes pluriclosed metrics and closed K\"ahler currents. These calculations supply geometric information about this example without giving a classification of all lower-mass degeneration.

The equation consequence is exact for a Hermitian background: \cref{cor:positive-lower-mass-solvability} characterizes positive lower mass by solvability for all finite non-pluripolar measures. When it vanishes, even a positive $L^1$ density bounded away from zero can admit no positive scaled subsolution. The remaining equation questions concern specified measures or additional geometric hypotheses, and separately the effect of the exponential weight.

\subsection{Primitive elliptic bundles}
\label{sec:primitive-elliptic-bundles}

On a primitive elliptic bundle, the variation of the top-degree mass is a vertical Dirichlet integral. This gives the exact upper mass and a positive lower bound for separated potentials. The lower mass over all potentials remains a geometric question, since the collapse construction over a product of curves does not transfer directly to a base without curves.

Let $(B,\eta)$ be a connected compact K\"ahler surface and let $L\to B$ be a holomorphic line bundle such that
\[
 \int_Bc_1(L)\wedge\eta=0.
\]
Choose a smooth Hermitian metric on $L$ whose curvature is primitive with respect to $\eta$. Such a metric exists: for an arbitrary metric, the trace of the curvature with respect to $\eta$ has mean zero by the condition above, and the scalar Poisson equation gives a smooth change of metric that cancels this trace. Let $c$ be the resulting smooth closed curvature representative of $-c_1(L)$, so that $\eta\wedge c=0$. Write $\pi\colon L\to B$ for the bundle projection and $w$ for a nonzero vector in a fiber, with norm taken in the chosen metric. On the complement $L^\times$ of the zero section, put
\[
 t=\log|w|^2,\qquad \ddc t=\pi^*c.
\]
For $q>1$, form the quotient
\[
 X=L^\times/\langle w\mapsto qw\rangle,
 \qquad \pi\colon X\longrightarrow B,
 \qquad T=2\log q,
 \qquad A=\int_B\eta^2,
\]
where we use the same symbol $\pi$ for the induced projection. The action increases $t$ by $T$. It is free and properly discontinuous, and the compact annulus $1\leq|w|\leq q$ meets every orbit. Hence $X$ is a connected compact complex threefold; in this subsection it replaces the general manifold of the chapter. Over a trivializing open set $U\subset B$, a local holomorphic frame of $L$ identifies the fibration with $U\times(\C^*/q^{\mathbb Z})$.
\Needspace{5\baselineskip}
Since $\mathrm{d}t$ and $\dc t$ are invariant under the action, they descend to $X$, and we set
\[
 D=\mathrm{d}t\wedge\dc t,\qquad \omega=\pi^*\eta+D.
\]
The form $D$ has rank one; it is semipositive and positive on the vertical line, and its kernel projects isomorphically onto the tangent space of $B$. Hence $\omega$ is Hermitian.

\Needspace{17\baselineskip}
\begin{proposition}\label{prop:primitive-elliptic-metric}\index{Elliptic fibration!primitive bundles}
The metric $\omega$ is balanced and satisfies
\[
 \ddc\omega=-\pi^*(c^2)\geq0,\qquad
 \int_X\omega^3=3TA.
\]
For every smooth real $u$,
\begin{equation}\label{eq:primitive-elliptic-mass}
 \int_X(\omega+\ddc u)^3
 =3TA-3\int_X\mathrm{d}u\wedge\dc u\wedge\bigl(-\pi^*(c^2)\bigr).
\end{equation}
In particular $v_+(\omega)=3TA$. The canonical bundles satisfy $K_X\simeq\pi^*K_B$. If $c_1(L)\ne0$ in real cohomology, $X$ admits neither a pluriclosed metric nor a closed K\"ahler current. If $c_1(L)=0$ in real cohomology, the chosen $\omega$ is K\"ahler.
\end{proposition}
\begin{proof}
Write $a=\mathrm{d}t$ and $b=\dc t$. Then $\mathrm{d}a=0$, $\mathrm{d}b=\pi^*c$, $\dc a=-\pi^*c$ and $\dc b=0$. Consequently
\[
 D^2=0,\qquad \mathrm{d}D=-a\wedge\pi^*c,
 \qquad \ddc D=-\pi^*(c^2),
\]
\[
 \mathrm{d}(\omega^2)=-2\pi^*(\eta\wedge c)\wedge a=0,
\]
so $\omega$ is balanced. Since $c$ is primitive, its eigenvalues relative to $\eta$ are $\lambda$ and $-\lambda$ for a function $\lambda\geq0$, so $c^2=-\lambda^2\eta^2$. This proves the asserted sign of $\ddc\omega$.

To compute the volume, fix an elliptic fiber $E=\pi^{-1}(b_0)$, where $b_0\in B$. In a logarithmic fiber coordinate $x+\mathrm{i}y$, the real periods are $\log q$ and $2\pi$, while $t$ equals $2x$ up to an additive constant and $D|_E=\pi^{-1}\mathrm{d}x\wedge\mathrm{d}y$. Thus $\int_ED=T$. Since $\omega^3=3\pi^*(\eta^2)\wedge D$, the volume formula follows.

For the mass identity, expand the cube in \eqref{eq:primitive-elliptic-mass}. The terms with one and with three factors $\ddc u$ integrate to zero, because $\omega^2$ is closed and $\ddc u$ is exact. For the remaining term, put $S=\ddc\omega$. Integration by parts gives
\[
 \int_X(\ddc u)^2\wedge\omega
 =\int_Xu\,\ddc u\wedge S
 =-\int_X\mathrm{d}u\wedge\dc u\wedge S,
\]
where the last equality uses $\mathrm{d}S=0$. This proves the identity. Its last integral is nonnegative, and $u=0$ attains the upper bound $3TA$. Since smooth and bounded potentials have the same mass extrema, by \cref{thm:gaud-volume-mass-extrema}, this proves $v_+(\omega)=3TA$.

If $c_1(L)=0$ in real cohomology, then $c$ is exact and Stokes' theorem gives $\int_Bc^2=0$. The pointwise identity $c^2=-\lambda^2\eta^2$ then forces $c=0$, and $\omega$ is closed. Otherwise $S=-\pi^*(c^2)$ is nonzero, smooth and semipositive. For a pluriclosed metric $\gamma$, Stokes' theorem would give
\[
 0<\int_X\gamma\wedge S
 =\int_X\omega\wedge\ddc\gamma=0,
\]
a contradiction. Likewise, a closed positive current $R\geq\varepsilon\omega$ with $\varepsilon>0$ would give $0<\int_XR\wedge S=\int_X\omega\wedge\ddc R=0$. All these pairings are pairings of currents with smooth forms.

Finally, let $w$ be a local holomorphic fiber coordinate given by a frame of $L$, and let $\sigma$ be a local holomorphic two-form on $B$. The form $\pi^*\sigma\wedge\mathrm{d}w/w$ is invariant under the quotient action. Changing the frame of $L$ adds to $\mathrm{d}w/w$ a horizontal one-form, whose wedge product with $\pi^*\sigma$ vanishes. These local sections of $K_X$ therefore have the transition functions of $\pi^*K_B$, which proves $K_X\simeq\pi^*K_B$.
\end{proof}

\Needspace{14\baselineskip}
For separated potentials, positivity bounds both the amplitude and the downward slope of the vertical derivative. The sharp variance bound under these two constraints determines how much mass such a potential can lose. For $M,T>0$, set
\[
 F(M,T)=
 \begin{cases}
  T^2/12,&T\leq2M,\\
  M^2-4M^3/(3T),&T\geq2M.
 \end{cases}
\]

\begin{lemma}\label{lem:periodic-variance-bound}
If a smooth real $T$-periodic function $p$ has mean zero, $|p|<M$ and $p'>-1$, then
\[
 \frac1T\int_0^T p^2\,\mathrm{d}s<F(M,T).
\]
The supremum of the left side over this class is $F(M,T)$.
\end{lemma}
\begin{proof}
For $p=0$ the inequality holds because $F(M,T)>0$. Otherwise set $a=\max p>0$, $-b=\min p<0$ and $r=a+b$. Integrating $-p'<1$ over a forward interval $I$ from a maximum to the next minimum gives $r<|I|\leq T$; also $r<2M$, since $|p|<M$. The function $G(s)=(a-s)(s+b)$ is nonnegative on the range of $p$. Since $p$ has mean zero and $-p'<1$, integration over $I$ gives
\[
 abT-\int_0^T p^2\,\mathrm{d}s
 =\int_0^T G(p)\,\mathrm{d}s
 \geq-\int_I G(p)p'\,\mathrm{d}s
 =\frac{r^3}{6}.
\]
As $ab\leq r^2/4$, the average of $p^2$ is at most $r^2/4-r^3/(6T)$. This expression increases strictly for $0<r<T$, and its value at $r=\min\{2M,T\}$ is $F(M,T)$. Since $r<\min\{2M,T\}$, the strict inequality follows.

For sharpness, consider the following $T$-periodic profiles $p_0$. When $T\geq2M$, let $p_0$ be constant equal to $M$ and to $-M$ on two intervals of length $(T-2M)/2$ each, joined by a descent of slope $-1$ and length $2M$, and close the period by an upward jump. When $T\leq2M$, let $p_0(s)=T/2-s$ for $0\leq s<T$, with its upward jump at the end of the period. In both cases $p_0$ has mean zero, $|p_0|\leq M$, $p_0'\geq-1$ in the sense of distributions, and the average of $p_0^2$ is $F(M,T)$; both constructions apply when $T=2M$. Convolution with a nonnegative periodic approximate identity preserves the mean and the bound on the distributional derivative, and multiplication by $1-\varepsilon$ makes both inequalities strict. As the convolution scale and $\varepsilon$ tend to zero, the resulting smooth functions converge to $p_0$ in $L^2$, which proves that the supremum is $F(M,T)$.
\end{proof}

\Needspace{18\baselineskip}
\begin{proposition}\label{prop:primitive-elliptic-separated}\index{Monge--Ampere mass@Monge--Amp\`ere mass!separated potentials}
Suppose $c_1(L)\ne0$ in real cohomology, and put
\[
 \kappa=-\frac{\int_Bc^2}{A}>0.
\]
Every smooth potential $u=f(t)+\pi^*v$ with $f$ of period $T$, $v\in C^\infty(B,\R)$ and $\omega+\ddc u>0$ satisfies
\[
 \int_X(\omega+\ddc u)^3
 >3TA\bigl(1-\kappa F(\kappa^{-1/2},T)\bigr)>0.
\]
Let $\lambda\geq0$ be the nonnegative eigenvalue function of $c$ relative to $\eta$, so $c^2=-\lambda^2\eta^2$, and put $M_0=1/\max_B\lambda$. For the subclass $v=0$, the exact mass infimum is
\[
 3TA\bigl(1-\kappa F(M_0,T)\bigr)>0.
\]
It is approached by smooth potentials $f(t)$ with $\omega+\ddc f(t)>0$, and it is not attained by such potentials.
\end{proposition}
\begin{proof}
Put $p=f'$. The chain rule gives
\[
 \omega+\ddc u
 =\pi^*(\eta+pc+\ddc v)+(1+p')D.
\]
Positivity on the vertical line implies $p'>-1$. Over every point of $B$, the function $t$ takes every value modulo $T$, so restriction to $\ker\partial t$ shows that $\eta+pc+\ddc v$ is a K\"ahler form on $B$ for every value of $p$. Its square has integral
\[
 \int_B(\eta+pc+\ddc v)^2=A(1-\kappa p^2)>0,
\]
since primitivity removes the mixed term and Stokes' theorem removes the exact terms. Hence $|p|<\kappa^{-1/2}$.

Only the vertical derivative enters \eqref{eq:primitive-elliptic-mass}, because $-\pi^*(c^2)$ has top degree in the base directions. For a $T$-periodic function $F_0$, integration over a fiber gives $\int_E F_0(t)D=\int_0^T F_0(s)\,\mathrm{d}s$, and therefore
\[
 \int_X(\omega+\ddc u)^3
 =3TA\left(1-\frac\kappa T\int_0^T p^2\,\mathrm{d}s\right).
\]
Now apply \cref{lem:periodic-variance-bound} with $M=\kappa^{-1/2}$. The formula for $F$ gives $F(M,T)<M^2$, which proves positivity of the bound.

For $v=0$, diagonalizing $c$ with respect to $\eta$ shows that the horizontal positivity condition is exactly $|p|<M_0$; together with $p'>-1$ it also suffices for positivity of the full form. Moreover $\kappa$ is the $\eta^2$-average of $\lambda^2$, so $\kappa\leq M_0^{-2}$. \Cref{lem:periodic-variance-bound} with $M=M_0$ therefore proves the claimed positive lower bound and its nonattainment. The smooth approximating functions in its proof have mean zero and hence periodic primitives $f_\varepsilon$, and the potentials $f_\varepsilon(t)$ prove that the bound is sharp.
\end{proof}

The proof of the first bound does not show that it is sharp for separated potentials over an arbitrary base: positivity of the classes $[\eta]+p[c]$ alone does not provide one function $v$ that makes all the representatives $\eta+pc+\ddc v$ positive. The threefold of \cref{sec:vanishing-lower-mass} is the primitive elliptic bundle over $\mathbb P^1\times\mathbb P^1$ with $q=\mathrm e^{2\pi}$, so that $A=2$, $\kappa=1$ and $T=4\pi$. For it, \cref{prop:primitive-elliptic-metric,prop:primitive-elliptic-separated} reduce to \cref{prop:lower-mass-balanced-metric,prop:lower-mass-no-pluriclosed,prop:elliptic-radial-mass}, the bound is $8$ and it is sharp, and by \cref{thm:vanishing-lower-mass} its lower mass vanishes for every Hermitian form.

A base without curves rules out the averages of horizontal curves used in the known collapse example. The following construction uses the period theorem for K3 surfaces and the description of their K\"ahler cone \cite[Theorems 7.4.1 and 8.5.2 of the online version]{HuyK3}.

\begin{example}\label{ex:primitive-elliptic-k3}\index{Elliptic fibration!primitive K3 example}
There is a K3 surface $B$ with $\operatorname{Pic}(B)=\mathbb Z[L]$, $c_1(L)^2=-4$, and a K\"ahler form $\eta$ orthogonal to $c_1(L)$. It has no compact curves. The associated $X$ has trivial canonical bundle, a balanced metric of finite upper mass, and no pluriclosed metric or closed K\"ahler current. It has no divisors, and every compact irreducible curve is an elliptic fiber.
\end{example}
\begin{proof}
In a hyperbolic plane of the K3 lattice, choose a basis $e,f$ with $e^2=f^2=0$ and $e\cdot f=1$, and put $\ell=e-2f$. This vector is primitive and has square $-4$. Choose a period in $\ell^\perp$ outside the hyperplanes perpendicular to integral vectors not in $\mathbb Z\ell$. Such a period exists by Baire's theorem: the period domain in $\ell^\perp$, of signature $(3,18)$, is nonempty and each excluded hyperplane is proper. The period theorem and the description of integral $(1,1)$-classes give the asserted Picard group. An irreducible curve would have nonzero class, because its K\"ahler area is positive, and hence square $-4m^2<-2$ for some integer $m\neq0$, contradicting adjunction. Thus there are no curves.

The orthogonal complement of $\ell$ in $H^{1,1}(B,\R)$ has signature $(1,18)$. Choose a class of positive square there, with sign chosen so that it lies in the component of the positive cone containing the K\"ahler classes. Since there are no rational curves, the K\"ahler-cone theorem shows that this class is K\"ahler; choose $\eta$ in it. \Cref{prop:primitive-elliptic-metric} then gives the assertions about the metric and the canonical bundle.

If a divisor $Z\subset X$ dominated $B$, the push-forward of its integration current would be $\pi_*[Z]=m$ for a positive generic degree $m$. Stokes' theorem and the projection formula would then give
\[
 0=\int_X[Z]\wedge\ddc\omega
   =m\int_B(-c^2)>0.
\]
Thus no divisor dominates $B$. The image of any other irreducible divisor would be a curve or a point; the base has no curves, and a fiber has dimension one, so both alternatives are impossible. Likewise, the image of an irreducible compact curve is a point, and the curve is therefore an entire elliptic fiber.
\end{proof}

Put $V=3TA$. The root volumes of complementary tests can be computed explicitly for one radial family. For smooth real $T$-periodic functions $a,g$, consider
\[
 \Gamma=a(t)\eta^2+g(t)\eta\wedge D.
\]

\Needspace{15\baselineskip}
\begin{proposition}\label{prop:primitive-radial-root-volume}\index{Gauduchon volume!primitive radial tests}
The form $\Gamma$ is strictly positive and $\ddc$-closed exactly when $a$ is a positive constant and $g>0$. For such a form put $\bar g=T^{-1}\int_0^Tg(s)\,\mathrm ds$ and $d_\Gamma=\int_X\omega\wedge\Gamma$. Its positive square root and volume are
\[
 \gamma=\sqrt a\,\eta+\frac{g(t)}{2\sqrt a}D,
 \qquad
 \int_X\gamma^3=\frac32TA\sqrt a\,\bar g
       \leq\frac{d_\Gamma^{3/2}}{\sqrt V}.
\]
Equality holds exactly when $\bar g=2a$, without requiring $g$ to be constant. In particular every degree-one test in this family has root volume at most $V^{-1/2}$.
\end{proposition}
\begin{proof}
Since $\eta\wedge c=0$, the form $g\eta\wedge D$ is closed. The chain rule gives
\[
 \ddc(a\eta^2)=a''D\wedge\eta^2,
\]
the $a'c\wedge\eta^2$ term having excess base degree. Thus $\ddc\Gamma=0$ forces $a''=0$, and periodicity makes $a$ constant. In a positive coframe with $\eta=\sigma_1+\sigma_2$, $D=\sigma_3$, the complementary coefficient matrix of $\Gamma$ is $\operatorname{diag}(g,g,2a)$. Hence positivity is exactly $a,g>0$.
The displayed $\gamma$ is positive and squares to $\Gamma$. Multiplication and fiber integration give
\[
 d_\Gamma=TA(a+\bar g),\qquad
 \gamma^3=\tfrac32\sqrt a\,g\eta^2\wedge D.
\]
For fixed $d_\Gamma>0$ set $s=d_\Gamma/(TA)$. Then $0<a<s$, $\bar g=s-a$, and $\sqrt a(s-a)$ has its unique maximum at $a=s/3$. Substitution proves the bound and its equality condition.
\end{proof}

For the divisor-free threefold of \cref{ex:primitive-elliptic-k3}, \cref{thm:divisor-free-transverse-concentration} gives divergent transverse pairings against normalized Gauduchon tests. The radial family above has bounded root volume at fixed degree, even when $g$ concentrates arbitrarily strongly. Large transverse pairings do not by themselves give unbounded root volumes or a sequence of potentials with masses tending to zero.

\begin{problem}\label{prob:primitive-elliptic-lower-mass}
Determine the lower Monge--Amp\`ere mass (\cref{def:bmp}) of the primitive elliptic bundles other than the threefold of \cref{sec:vanishing-lower-mass}, for instance of their balanced metrics $\omega$ from \cref{prop:primitive-elliptic-metric}, in particular for the bundles of \cref{ex:primitive-elliptic-k3}. Does it vanish whenever $c_1(L)\neq0$ in real cohomology?
\end{problem}

The estimates for separated potentials do not answer this question; the known collapse over $\mathbb P^1\times\mathbb P^1$ uses potentials outside that class. On the K3 example every compact curve is vertical, so averaging over such curves cannot produce the horizontal directions used in \cref{lem:lower-mass-curve-average}. A proof of collapse would therefore need a different construction, while positive lower mass would require a bound for all admissible potentials or for all normalized Gauduchon tests.


\subsection{Null loci and the singularities of the Chern--Ricci flow}

Let $X$ be a connected compact complex manifold of dimension $n$, let $\theta$ be a smooth real $(1,1)$-background, and fix a smooth strictly positive probability form $\mu$. Write $\mathcal V_\mu$ for the volume in \cref{def:gaud-variational-volume}. In discussing the flow, $\omega_0$ is the initial Hermitian metric, $\operatorname{Ric}(\omega_0)$ its conventional Chern--Ricci form, and $\tau$ the finite maximal time; thus $\theta_\tau=\omega_0-\tau\operatorname{Ric}(\omega_0)$. An annihilating current $T$ means a positive $\ddc$-closed $(n-1,n-1)$-current normalized by the fixed Hermitian trace in \cref{cor:chern-ricci-current-lifetime}.

For a big smooth background $\theta$, Dang defines its non-Hermitian locus $E_{\mathrm{nH}}(\theta)$ as the intersection of the singular sets of all strict currents in that background with \emph{almost analytic singularities}; this intersection is itself attained by such a current. In the source's terminology the potential is locally a positive multiple of the logarithm of an ideal norm, up to a bounded remainder, and on a suitable log resolution its remainder after subtracting the divisorial logarithms is smooth. We retain this exact singularity condition.

\Needspace{8\baselineskip}
For nef and big $\theta$, define
\[
 \operatorname{Null}(\theta)
 =\bigcup_{\substack{V\subset X\text{ irreducible analytic}\\
                     \dim V>0,\ \theta|_V\text{ not big}}}V.
\]
Bigness on a singular subvariety uses its intrinsic analytic definition. Dang proves $E_{\mathrm{nH}}(\theta)=\operatorname{Null}(\theta)$ for arbitrary compact complex $X$ and nef, big, possibly nonclosed $\theta$. This qualitative identification is therefore a known theorem. Detecting bigness and this locus from numerical data is a further issue: the numerical criterion recorded in \cref{sec:history-geometry} uses bounded mass on $X$ and positive lower volume on its positive-dimensional submanifolds.

The known flow theorem recorded in \cref{sec:history-geometry} identifies finite-time Chern--Ricci singularities with the null locus under positive limiting lower volume, for arbitrary initial metrics in dimension two, and for Guan--Li initial metrics in higher dimensions. Here Guan--Li means $\ddc\omega_0=0$ and $\mathrm{d}\omega_0\wedge \dc\omega_0=0$. These hypotheses must not be silently replaced by positivity of $\mathcal V_\mu$.

The current calculation in \cref{cor:chern-ricci-current-lifetime} applies without those extra hypotheses. At a finite maximal time $\tau$, the background $\theta_\tau=\omega_0-\tau\operatorname{Ric}(\omega_0)$ is nef. Its normalized annihilating currents satisfy
\[
 \int_X\theta_\tau\wedge T=0.
\]
If $\mathcal V_\mu(\theta_\tau)>0$, their determinant-weighted measures vanish by \cref{prop:nef-annihilator-determinant}. This is a precise degeneracy statement, but it does not imply that their traces are concentrated on analytic sets. A smooth positive semidefinite matrix of deficient rank can be defined on an open set.

\Needspace{8\baselineskip}
\begin{problem}\label{qu:chern-ricci-gaud-singularities}
Suppose a Chern--Ricci flow has finite maximal time $\tau$ and $\mathcal V_\mu(\theta_\tau)>0$. Can one prove that $\theta_\tau$ is big and that the flow converges smoothly to a positive metric locally on $X\setminus E_{\mathrm{nH}}(\theta_\tau)$, without a bounded mass hypothesis? How do the annihilating Gauduchon currents describe the resulting degeneracy?
\end{problem}

The first assertion is a case of the volume-to-bigness question. Once bigness of $\theta_\tau$ is known, the parabolic theorem recorded in \cref{sec:history-geometry} supplies convergence as currents and smooth local convergence off $E_{\mathrm{nH}}(\theta_\tau)$, by parabolic estimates. Combined with nefness and the qualitative null-locus theorem, this identifies the analytic exceptional set. Thus the missing step in this route is the proposed numerical criterion for bigness. The local envelope continuity and homogeneous replacement arguments of \cref{thm:big-continuous-envelope-contact} supply an elliptic foundation away from a strict barrier's poles; they are not a reconstruction of those parabolic estimates.

Completeness requires an additional metric estimate. In \cref{thm:incomplete-divisorial-metric}, the unique finite-energy solution of a divisorial exponential equation in a big closed class is incomplete on every nonempty connected open subset of its ample locus. Its minimal divisorial singularity cancels the prescribed pole. This resolves the unrestricted big-class completeness question negatively; it does not decide completeness under nefness or under the compatibility conditions of the geometric smoothing problem discussed in \cref{sec:history-completeness}. In particular it does not settle \cref{qu:chern-ricci-gaud-singularities}, whose limiting background is nef.

For integrable poles weakened by logarithmic factors, \cref{cor:divisorial-exponential-completeness} gives a sufficient condition in terms of the density, the potential and the tangential determinant. The proof controls normal length and applies to arbitrary Hermitian backgrounds. The torus models in \cref{ex:critical-complete-divisor,ex:supercritical-incomplete-divisor} attain its critical exponent and show its sharpness, even with a K\"ahler background and full mass. Their pole coefficients need not extend smoothly and positively across the divisor; this distinction preserves the separate scope of the geometric smoothing question.

\subsection{Degenerate densities and uniqueness}

Let $X$ be a connected compact complex manifold, let $\omega$ be a Hermitian form and let $\mathrm dV$ be a smooth positive volume form. For $u\in\PSH(X,\omega)$ write $\mathrm{MA}_\omega(u)$ for its Monge--Amp\`ere measure. Densities $f$ below are smooth nonnegative functions on $X$; normalization of a potential means $\sup_Xu=0$. The scalar $c$ and exponential parameter $\lambda$ are positive real numbers.

The counterexample in \cref{thm:degenerate-hermitian-nonuniqueness} shows that a smooth nonnegative density can have two distinct smooth normalized solutions even on $\mathbb P^3$. The further construction in \cref{thm:numa-thin-zero-set} has uncountably many smooth normalized solutions and a density positive off two real hypersurfaces. By modifying the background, \cref{thm:quadratic-density-nonuniqueness} achieves exactly quadratic vanishing on those hypersurfaces. Thus neither zero smooth volume of the zero set nor finite-order vanishing suffices for uniqueness. All the backgrounds are strictly Hermitian and have bounded mass; their non-Hermitian loci are empty. The degeneration causing nonuniqueness belongs to the density and the solution forms, and is not detected by the background's null locus.

This places a necessary restriction on any geometric use of the equation. A construction that assigns a canonical potential to arbitrary nonnegative data must specify a selection procedure and prove that it is well defined. Existence in the ordinary minimal singularity type does not provide such a selection. The unique scalar, and the corresponding variational volume for strictly positive data, remain meaningful without uniqueness of every potential. On a compact Hermitian surface, \cref{uni-cor:dense-lowdim} proves uniqueness for $C^2$ potentials and continuous density positive on a dense set, with no differential condition on the metric. In dimension three, \cref{uni-thm:sobolev} instead uses a sign condition on $\ddc\omega$. The cofactor identity in \eqref{eq:mixed-uniqueness-defect} and the signed quadratic-difference argument provide two distinct integration mechanisms.

\begin{problem}\label{qu:density-zero-set-uniqueness}
Let $X$ be connected, $\omega$ Hermitian and $f\geq0$ smooth. Suppose $\{f>0\}$ is connected and dense. Are two bounded normalized solutions of $\mathrm{MA}_\omega(u)=cf\,\mathrm{d}V$, with the same $c>0$, necessarily equal? Separately, for densities allowing nonuniqueness, do the solutions of $\mathrm{MA}_\omega(u_\lambda)=\mathrm{e}^{\lambda u_\lambda}f\,\mathrm{d}V$ converge uniformly after subtracting their suprema, as $\lambda\downarrow0$?
\end{problem}

For $C^{1,1}$ potentials the first question is answered affirmatively by \cref{uni-thm:anchors}. The remaining issue is the passage to bounded potentials. All the hypersurface counterexamples here have three positive-density components; even the analytic density in \cref{uni-thm:analytic-example} has exactly this topology. The second asks for a canonical limit, and does not follow from uniqueness of the positive-exponent equation at each fixed parameter. A selection could exist even when the limiting equation has several solutions. These counterexamples do not refute \cref{conj:continuous-gaud-adjoint}: that conjecture uses a finite continuous logarithmic density $F$, for which $\mathrm{e}^F$ is strictly positive. The logarithms of the degenerate densities have no finite continuous extensions across their zero sets.

There is already some control of selection in the Hermitian $L^p$ setting: \cref{prop:symmetric-equation-clusters} proves uniform precompactness of the normalized approximations, and every cluster limit preserves the symmetries of the data. For the symmetric projective example, the radial argument of \cref{thm:numa-selected-limit} goes further: the whole normalized family converges uniformly to the explicit symmetric $C^{1,1}$ solution of \cref{prop:numa-radial-continuum}. Thus this example has both many normalized solutions and a definite limit for the specified approximation. The general selection question above remains open in this treatment.

For a nonzero density at positive threshold in dimension at least two, \cref{cor:smooth-dual-exact-equation,cor:smooth-dual-exponential-limit} give uniqueness of the normalized threshold solution and uniform convergence of the whole normalized positive-exponent family whenever a smooth strictly positive maximizing Gauduchon current exists. The construction in \cref{ex:nonclosed-dual-zero-density} satisfies this condition even though the density vanishes on an open set. Thus the size of the zero set alone does not determine uniqueness; the dual condition also involves the background.

\subsection{Balanced and pluriclosed metrics}

Let $X$ be a connected compact complex manifold of dimension $n\geq2$. The forms $\beta$ and $\omega$ below are Hermitian metrics on $X$, and $c_1(X)$ denotes the first real Chern class of its holomorphic tangent bundle. The Fujiki class $\mathcal C$ consists of manifolds bimeromorphic to compact K\"ahler manifolds.

A Hermitian metric $\beta$ is balanced if $\mathrm{d}\beta^{n-1}=0$, while a Hermitian metric $\omega$ is pluriclosed if $\ddc\omega=0$. The Fino--Vezzoni conjecture asks whether a compact complex manifold admitting both types of metrics is K\"ahler; the two metrics can be different. Its source is recorded in \cref{sec:history-geometry}. A single metric satisfying both conditions is already K\"ahler, which does not resolve the question with two metrics.

Chiose proves that a compact manifold in $\mathcal C$ admitting a pluriclosed metric is K\"ahler. Consequently, in the balanced--pluriclosed setting, it would suffice to construct a closed big $(1,1)$-current and apply the characterization of $\mathcal C$. This route emphasizes a geometric existence problem which is absent when one only works with an arbitrary Hermitian background.

There is an elementary obstruction to using the scalar equation alone. For every distributional potential $u$,
\[
 \mathrm{d}(\omega+\ddc u)=\mathrm{d}\omega.
\]
Thus if the starting pluriclosed metric is nonclosed, no potential in its fixed background makes it K\"ahler. The finite-volume equation, including its smooth solution, does not change this differential datum. Similarly, a conformal Gauduchon representative need not be balanced. Positivity and a scalar normalization cannot be used to replace either missing closedness condition.

The balanced form $\beta^{n-1}$ nevertheless supplies more than a Gauduchon test: it pairs trivially with every exact real two-form by Stokes. It therefore interacts with flows whose variation has an exact component. Fino--Vezzoni propose long-time existence and convergence to a K\"ahler metric for the pluriclosed flow when a balanced metric exists and $c_1(X)=0$. Their invariant examples illustrate the use of a conserved balanced trace. In general that integral is not a pointwise metric bound.

\begin{problem}\label{qu:balanced-pluriclosed-program}
Use the coexistence of a balanced metric and a pluriclosed metric to construct closed positive data, or to control a geometric flow that can change the torsion. Determine what estimate beyond the finite balanced trace prevents metric degeneration. Can determinant estimates for Gauduchon currents contribute to that estimate?
\end{problem}

This asks for an interaction with the scalar theory, not an implication already proved by its equation. One can first study families for which the balanced trace and the determinant density are explicitly controlled. The passage from an integral trace bound to local ellipticity is then a concrete analytic target. No such passage for arbitrary families is asserted here.

\subsection{Morse inequalities and the choice of volume}

For the closed comparison, let $X$ be a connected compact K\"ahler $n$-fold and let $\alpha,\beta$ be closed nef real $(1,1)$-classes. Write $\operatorname{Vol}$ for the usual volume of a closed class. In the Hermitian discussion, $X$ is only assumed connected, compact and complex; $\theta,\eta$ are smooth real $(1,1)$-forms, and $\mu$ is a fixed smooth strictly positive probability form. Use $\lambda_\mu$, $\mathcal V_\mu$ and $D_\mu$ from \cref{def:gaud-variational-volume,sec:gauduchon-collapse}.

The closed K\"ahler theory provides a demanding comparison. Tosatti's recent preprint establishes the transcendental Morse inequality
\[
 \operatorname{Vol}(\alpha-\beta)
 \geq\int_X\alpha^n-n\int_X\alpha^{n-1}\wedge\beta
\]
for closed nef classes on compact K\"ahler manifolds, together with the resulting volume differentiability and cone-duality consequences. We record it as a recent literature result; its proof is not reconstructed here. The unrestricted non-K\"ahler extension remains a different problem. The qualitative bounded-mass criterion and the source of the closed theorem are recorded in \cref{sec:history-geometry}.

For nonclosed backgrounds, the displayed intersection expression is not generally invariant under changing representatives. Also, the upper class volume and $\mathcal V_\mu$ have different definitions and different possible finiteness behavior. A Hermitian analogue must first specify which invariant it estimates.

\Needspace{10\baselineskip}
The first-variation formula gives a finite quantity to investigate. Fix a real number $t>0$. Suppose $\theta_s=\theta-s\eta$ has a smooth positive representative for $0\leq s\leq t$, and let $\gamma_s$ be a minimizing Gauduchon metric for its smooth equation. By \cref{prop:gaud-volume-first-variation} and the fundamental theorem of calculus,
\begin{equation}\label{eq:gaud-volume-subtraction-cost}
 \lambda_\mu(\theta-t\eta)
 =\lambda_\mu(\theta)
   -\int_0^t
       \frac{\int_X\eta\wedge\gamma_s^{n-1}}
            {D_\mu(\gamma_s)}\,\mathrm{d}s.
\end{equation}
The quotient is unchanged by rescaling the minimizer. This identity is valid on the stated positive segment and uses no upper Monge--Amp\`ere mass bound. A uniform upper bound on its integrand gives a lower bound for the remaining root volume. Proving a useful bound near a degenerate endpoint is an additional problem; concavity alone provides the tangent upper bound, not the lower bound needed here.

\begin{problem}\label{qu:hermitian-morse-program}
Find geometric hypotheses controlling the cost in \eqref{eq:gaud-volume-subtraction-cost} up to a nef endpoint, and determine whether the resulting positive root volume forces bigness there. For closed backgrounds, compare the resulting estimate with classical Morse inequalities. For nonclosed backgrounds, retain the differential data and state explicitly any dependence on $\mu$.
\end{problem}

The singularity and divisorial constructions in the book become relevant at this stage: they record how much positivity can be spent on imposed poles, while the residual Gauduchon volume stays finite. They do not yet give the moving intersection product or the birational volume formula of the closed theory. Recovering such geometric information is a further objective, with the finite-divisor and Hopf calculations serving as explicit tests.
